\documentclass[conference]{IEEEtran}
% Uncomment the line below if targeting a full-length IEEE Transactions journal submission:
% \documentclass[journal]{IEEEtran}

\IEEEoverridecommandlockouts

% Standard IEEE and scientific packages
\usepackage{cite}
\usepackage{amsmath,amssymb,amsfonts}
\usepackage{algorithmic}
\usepackage{graphicx}
\usepackage{textcomp}
\usepackage{xcolor}
\usepackage{booktabs}
\usepackage{multirow}
\usepackage{tabularx}
\usepackage{url}
\usepackage[capitalize]{cleveref}

% Search path for figures provided to PaperOrchestra
\graphicspath{{figures/}}

\begin{document}

\title{Multi-Domain Transient and Grid-Forming Control Benchmark on the WSCC 9-Bus System}

% Anonymized author block for double-blind review (replace with author details for camera-ready):
\author{\IEEEauthorblockN{Anonymous Authors}
\IEEEauthorblockA{\textit{Department of Electrical Engineering} \\
\textit{Institution Name}\\
City, Country \\
email@domain.edu}}

\maketitle

\begin{abstract}
% To be populated by the Section Writing Agent from idea.md and experimental_log.md.
\end{abstract}

\begin{IEEEkeywords}
% To be populated with 4--6 IEEE index terms (e.g., Grid-forming inverter, virtual synchronous generator, synchronverter, weak grid, electromagnetic transient simulation, WSCC 9-bus system).
\end{IEEEkeywords}

\section{Introduction}\label{sec:intro}
% To be populated by the Literature Review Agent (macro-level context, problem gap, and core contributions).

\section{Related Work}\label{sec:related}
% To be populated by the Literature Review Agent (micro-level comparison across PLL weak-grid stability, RoCoF-VSG, synchronverters, and multi-domain EMT/phasor benchmarks).

\section{Multi-Domain WSCC 9-Bus System and Control Formulations}\label{sec:methodology}
% To be populated by the Section Writing Agent (hybrid EMT/phasor architecture, Models 0--3 equations, tuned parameters, and contingency definitions S1--S4).

\section{Cross-Scenario Transient and Power-Quality Evaluation}\label{sec:experiments}
% To be populated by the Section Writing Agent (10-KPI cross-scenario table, S1--S4 transient dynamics, weak-grid GFL collapse, and IEEE 519 harmonic compliance).

\section{Decoupled $(H, D)$ Parametric Sensitivity Analysis}\label{sec:sensitivity}
% To be populated by the Section Writing Agent (32-point parametric sweep isolating virtual inertia H and damping D_pu).

\section{Conclusion}\label{sec:conclusion}
% To be populated by the Section Writing Agent (summary of findings and topology selection guidelines).

\bibliographystyle{IEEEtran}
\bibliography{references}

\end{document}